\documentclass[11pt,a4paper]{article}
\usepackage[top=0.6in, bottom=0.6in, left=0.75in, right=0.75in]{geometry}
\usepackage[T1]{fontenc}
\usepackage[utf8]{inputenc}
\usepackage{enumitem}
\usepackage{hyperref}
\usepackage{microtype}
\usepackage{xcolor}
\usepackage{titlesec}

\definecolor{accent}{RGB}{0,100,80}

\titleformat{\section}{\large\bfseries\color{accent}}{}{0em}{}[\color{accent}\rule{\textwidth}{0.5pt}]
\titlespacing{\section}{0pt}{8pt}{4pt}

\setlist{nosep, leftmargin=*}
\pagestyle{empty}
\hypersetup{colorlinks=true, urlcolor=accent}

\begin{document}

\begin{center}
  {\LARGE \textbf{Ethan Brown, Ph.D.}}\\[3pt]
  {\large \textcolor{accent}{Research Scientist -- Biotech}}\\[4pt]
  Cambridge, MA \quad|\quad
  \href{mailto:ethan.brown@email.com}{ethan.brown@email.com} \quad|\quad
  (617) 555-0372 \quad|\quad
  \href{https://linkedin.com/in/ethanbrown-biotech}{LinkedIn}
\end{center}

\section{Summary}
Translational research scientist with Ph.D. in molecular biology and 6 years of industry experience in gene therapy and antibody drug development. Expertise in CRISPR screening, AAV manufacturing, flow cytometry, and IND-enabling studies. Co-inventor on 2 patents.

\section{Experience}

\noindent\textbf{Senior Research Scientist -- Gene Therapy} \hfill \textit{Feb 2021 -- Present}\\
\textit{Editas Medicine, Cambridge, MA}
\begin{itemize}
  \item Lead in vivo research program for CRISPR base editing therapy targeting CEP290 (LCA10); designed and executed mouse and non-human primate efficacy and toxicology studies.
  \item Developed high-throughput AAV production process increasing yield by 3x while reducing COGS by 40\%; transferred process to GMP manufacturing.
  \item Optimized LNP formulation for liver-targeted mRNA delivery achieving 85\% hepatocyte transfection in vivo; results supported IND filing for LCA10 program (2024).
  \item Supervised 2 scientists and 1 research associate; mentored associate scientist through independent research project resulting in co-authorship on \textit{Nature Medicine} paper.
\end{itemize}

\vspace{4pt}
\noindent\textbf{Research Scientist} \hfill \textit{July 2017 -- Jan 2021}\\
\textit{AstraZeneca, Cambridge, MA}
\begin{itemize}
  \item Generated and characterized 80+ antibody candidates against validated oncology targets; advanced 4 candidates through lead optimization to bispecific antibody stage.
  \item Established cell line development platform for antibody production (CHO, HEK293); reduced development timeline from 12 to 7 weeks.
  \item Contributed to 2 patent filings for novel antibody-antigen combinations.
\end{itemize}

\section{Technical Skills}
\textbf{Molecular Biology:} CRISPR-Cas9/base editing, cloning, lentiviral/AAV vector production, qPCR, WB, ELISA\\
\textbf{Cell Biology:} Primary cell culture (iPSC, T cells, hepatocytes), flow cytometry (15-color panels), live imaging\\
\textbf{In Vivo:} Mouse (C57BL/6, transgenic), NHP study management, retinal/systemic delivery, necropsy\\
\textbf{Analytics:} FlowJo, GraphPad Prism, R (basic), ImageJ, SnapGene

\section{Publications}
\begin{enumerate}[leftmargin=*, label={[\arabic*]}]
  \item \textbf{E. Brown}, K. Maeder, et al. ``Efficient Base Editing in the Retina via Subretinal AAV Delivery.'' \textit{Nature Medicine} 30(1), 2024.
  \item A. Liu, \textbf{E. Brown}, et al. ``Optimized LNP Formulations for Hepatocyte-Targeted mRNA Delivery.'' \textit{Molecular Therapy} 29(8), 2021.
\end{enumerate}

\section{Education}
\textbf{Ph.D. Molecular Biology}, MIT \hfill \textit{May 2017}\\
Dissertation: \textit{CRISPR-Mediated Correction of Retinal Dystrophy in Mouse Models}\\
Advisor: Prof. Feng Zhang \quad|\quad NSF Graduate Research Fellow\\[3pt]
\textbf{B.S. Biochemistry}, Johns Hopkins University \hfill \textit{May 2012} \quad GPA: 3.92/4.0

\end{document}
